% !TEX root = ../manuscript.tex

\section{Results}

The analysis is organized around the reservoir-scale questions enabled by the
framework. We first identify geologic controls on typical migration and
containment, then quantify realization-to-realization variability within each
geology. Targeted enrichment resolves conditional outcome distributions in
geologies with mixed behavior, and the deterministic benchmarks test when a
representative fault preserves the ensemble-level conclusion.

\subsection{Geologic controls on typical migration and containment}

\todo{Lead with one primary containment quantity and one primary migration
quantity. Report design-based main effects and selected interactions for
thickness structure, faulting depth, sand-layer clay volume fraction, and
clay-layer clay volume fraction. For zero-inflated quantities, separate the
occurrence of fault or top-seal entry from the magnitude after entry. Report
standardized effects and uncertainty intervals, and describe these as effects
over the prescribed, unweighted geologic design space rather than as
field-calibrated probabilities.}

\subsection{Within-geology variability from stochastic fault-core structure}

\todo{For every geology, report robust conditional spread, outcome-class
frequencies, and whether the 12 realizations cross a physically defined
containment threshold. Classify geologies using predefined categories such as
consistently contained, mixed containment outcomes, and consistently
migrating. Support this classification with the balanced design-based
decomposition of between-geology and average within-geology variability.}

\subsection{Conditional outcome distributions in enriched geologies}

\todo{Describe the prespecified enrichment criteria and report empirical CDFs,
quantiles with bootstrap intervals, and outcome frequencies with binomial
intervals for the selected geologies. Use newly added realizations to confirm,
rather than merely redisplay, realization sensitivity first observed in Phase
1. Include consistently contained and consistently migrating controls so the
selection cannot be interpreted as focusing only on anomalous cases.}

\subsection{Performance of representative-fault benchmarks}

\todo{Compare each full-distribution-medoid result with the conditional
stochastic distribution. Report its difference from the stochastic median,
its empirical CDF position, and its agreement with the majority outcome class.
Treat the low/high fields as unweighted stress benchmarks rather than bounds,
percentiles, or probabilistic realizations. Emphasize the conditions under
which a representative fault is adequate and those under which it gives a
misleading single answer.}
